% Mathématiques 4e — cours V3
% Nombres rationnels et calcul

\section{Objectifs}

\begin{itemize}
\item Étendre les quatre opérations aux nombres relatifs.
\item Déterminer le signe d'un produit ou d'un quotient.
\item Comprendre qu'un nombre rationnel est un quotient d'entiers.
\item Utiliser l'inverse d'un nombre non nul.
\item Multiplier et diviser des fractions.
\item Comparer des nombres rationnels écrits sous différentes formes.
\item Simplifier une fraction et reconnaître une fraction irréductible.
\item Résoudre des problèmes faisant intervenir plusieurs opérations sur des rationnels.
\end{itemize}

\section{1. Du nombre décimal au nombre rationnel}

Tous les quotients de deux entiers ne possèdent pas une écriture décimale finie.

Par exemple :
\[
1\div3=0{,}333\ldots
\]

Le nombre est pourtant parfaitement défini par :
\[
\dfrac13.
\]

\begin{definition}
Un nombre rationnel est un nombre qui peut s'écrire :
\[
\dfrac ab
\]
où $a$ est un entier relatif et $b$ un entier relatif non nul.
\end{definition}

Ainsi :
\[
-2{,}5=-\dfrac52
\]
est rationnel, tout comme :
\[
\dfrac7{12}.
\]

\section{2. Produit de deux nombres relatifs}

La valeur absolue du produit se calcule comme avec des nombres positifs.

Le signe dépend des signes des facteurs.

\coursefigure{/images/mathematiques/4eme/rationnels-signes.svg}{Tableau des signes d'un produit ou d'un quotient de deux nombres relatifs.}{Deux nombres de même signe donnent un résultat positif ; deux nombres de signes différents donnent un résultat négatif.}

\begin{propriete}
Pour un produit de deux nombres relatifs :
\begin{itemize}
\item mêmes signes : produit positif ;
\item signes différents : produit négatif.
\end{itemize}
\end{propriete}

Exemples :
\[
(-7)\times(-4)=28,
\]
\[
(-7)\times4=-28.
\]

\section{3. Produit de plusieurs facteurs}

Pour plusieurs facteurs non nuls, on peut compter le nombre de facteurs négatifs.

\begin{propriete}
Un produit est positif si le nombre de facteurs négatifs est pair.

Il est négatif si le nombre de facteurs négatifs est impair.
\end{propriete}

Par exemple :
\[
(-2)\times(-3)\times5=30.
\]

Il y a deux facteurs négatifs.

Au contraire :
\[
(-2)\times(-3)\times(-5)=-30.
\]

\section{4. Quotient de nombres relatifs}

Les mêmes règles de signe s'appliquent au quotient.

\[
\dfrac{-24}{6}=-4,
\]
\[
\dfrac{-24}{-6}=4.
\]

\begin{attention}
Le quotient par zéro n'existe pas.
\end{attention}

On ne peut jamais écrire :
\[
\dfrac a0.
\]

\section{5. Opposé et signe d'une fraction}

Une fraction peut porter son signe à plusieurs endroits :
\[
-\dfrac35
=
\dfrac{-3}{5}
=
\dfrac3{-5}.
\]

Mais :
\[
\dfrac{-3}{-5}
=
\dfrac35.
\]

Deux signes négatifs se compensent dans un quotient.

\section{6. Inverse}

\begin{definition}
L'inverse d'un nombre non nul $a$ est :
\[
\dfrac1a.
\]

Il est caractérisé par :
\[
a\times\dfrac1a=1.
\]
\end{definition}

L'inverse de :
\[
5
\]
est :
\[
\dfrac15.
\]

L'inverse de :
\[
-\dfrac37
\]
est :
\[
-\dfrac73.
\]

\section{7. Diviser revient à multiplier par l'inverse}

Pour $b\neq0$ :
\[
a\div b=a\times\dfrac1b.
\]

Cette propriété permet de définir la division des fractions.

\coursefigure{/images/mathematiques/4eme/rationnels-division-fractions.svg}{Transformation d'une division de fractions en multiplication par l'inverse.}{Diviser par une fraction non nulle revient à multiplier par la fraction renversée.}

\section{8. Multiplication de fractions}

Pour $b\neq0$ et $d\neq0$ :
\[
\boxed{
\dfrac ab\times\dfrac cd=\dfrac{ac}{bd}
}.
\]

\begin{exemple}
\[
\dfrac34\times\dfrac{10}{9}
=
\dfrac{30}{36}
=
\dfrac56.
\]
\end{exemple}

On peut parfois simplifier avant de multiplier :
\[
\dfrac{14}{15}\times\dfrac{25}{21}.
\]

On simplifie $14$ avec $21$ par $7$, et $25$ avec $15$ par $5$ :
\[
\dfrac23\times\dfrac53=\dfrac{10}{9}.
\]

\section{9. Division de fractions}

Pour $c\neq0$ :
\[
\boxed{
\dfrac ab\div\dfrac cd
=
\dfrac ab\times\dfrac dc
}.
\]

\begin{exemple}
\[
\dfrac56\div\dfrac{10}{9}
=
\dfrac56\times\dfrac9{10}
=
\dfrac{45}{60}
=
\dfrac34.
\]
\end{exemple}

\section{10. Addition et soustraction}

Les règles vues en 5e restent valables.

Pour additionner ou soustraire des fractions, il faut disposer d'un dénominateur commun.

\[
-\dfrac23+\dfrac54
=
-\dfrac8{12}+\dfrac{15}{12}
=
\dfrac7{12}.
\]

Le signe du numérateur fait partie du calcul.

\section{11. Priorités opératoires}

Dans une expression contenant plusieurs opérations :
\begin{enumerate}
\item on traite les parenthèses ;
\item puis multiplications et divisions ;
\item enfin additions et soustractions.
\end{enumerate}

Par exemple :
\[
A=-\dfrac35\times\dfrac{10}{9}-\dfrac13.
\]

On calcule d'abord le produit :
\[
-\dfrac35\times\dfrac{10}{9}
=
-\dfrac{30}{45}
=
-\dfrac23.
\]

Puis :
\[
A=-\dfrac23-\dfrac13=-1.
\]

\section{12. Comparer des nombres rationnels}

On peut utiliser plusieurs stratégies :
\begin{itemize}
\item même dénominateur ;
\item comparaison à zéro ou à un ;
\item écriture décimale lorsque celle-ci est pratique ;
\item mise au même dénominateur.
\end{itemize}

Par exemple :
\[
\dfrac7{12}
\]
et :
\[
\dfrac58.
\]

On utilise le dénominateur commun $24$ :
\[
\dfrac7{12}=\dfrac{14}{24},
\]
\[
\dfrac58=\dfrac{15}{24}.
\]

Donc :
\[
\dfrac7{12}<\dfrac58.
\]

\section{13. Fraction irréductible}

\begin{definition}
Une fraction est irréductible lorsque son numérateur et son dénominateur n'ont aucun diviseur commun autre que $1$.
\end{definition}

Par exemple :
\[
\dfrac{18}{24}
\]
n'est pas irréductible.

On simplifie par $6$ :
\[
\dfrac{18}{24}=\dfrac34.
\]

\section{14. Différentes représentations}

Un même rationnel peut être représenté par une fraction, une écriture décimale, un pourcentage ou un point sur une droite graduée.

Par exemple :
\[
\dfrac38=0{,}375=37{,}5\%.
\]

Choisir la représentation adaptée peut simplifier un problème.

\section{15. Problème de recette}

Une recette utilise :
\[
\dfrac34\ \text{L}
\]
de lait.

On prépare :
\[
\dfrac23
\]
de la recette.

Le volume nécessaire est :
\[
\dfrac34\times\dfrac23
=
\dfrac12.
\]

Il faut :
\[
\boxed{\dfrac12\ \text{L}}.
\]

\section{16. Problème de calcul en plusieurs étapes}

Calculons :
\[
B=\dfrac54-\left(-\dfrac23\right)\times\dfrac98.
\]

La multiplication est prioritaire :
\[
-\dfrac23\times\dfrac98=-\dfrac34.
\]

Donc :
\[
B=\dfrac54+\dfrac34=\dfrac84=2.
\]

\section{17. Contrôler un résultat}

Pour :
\[
\dfrac45\times\dfrac23,
\]
les deux facteurs sont compris entre $0$ et $1$.

Le produit doit donc être inférieur à chacun d'eux.

Le résultat :
\[
\dfrac8{15}
\]
est cohérent.

\section{18. Justifier une propriété}

Les règles de fractions ne sont pas de simples recettes.

Par exemple, si $c\neq0$ :
\[
\dfrac{ac}{bc}=\dfrac ab.
\]

On peut le justifier à partir de la définition du quotient.

\begin{remarque}
Vérifier une propriété sur plusieurs exemples ne constitue pas une démonstration générale.
\end{remarque}

\section{19. Méthode pour un calcul complexe}

\begin{methode}
\begin{enumerate}
\item repérer les opérations prioritaires ;
\item déterminer les signes ;
\item simplifier les fractions si possible ;
\item effectuer produits ou quotients ;
\item traiter additions et soustractions ;
\item rendre la fraction irréductible si demandé ;
\item contrôler le signe et l'ordre de grandeur.
\end{enumerate}
\end{methode}

\section{20. Erreurs fréquentes}

\begin{itemize}
\item Appliquer une règle de signe à une addition.
\item Oublier que deux signes négatifs donnent un quotient positif.
\item Inverser la mauvaise fraction lors d'une division.
\item Diviser par zéro.
\item Additionner les dénominateurs.
\item Simplifier à travers une addition.
\item Donner une fraction réductible lorsqu'une forme irréductible est demandée.
\end{itemize}

\section{À retenir}

\begin{aretenir}
Pour un produit ou quotient :
\[
\boxed{\text{mêmes signes}\Rightarrow +,\qquad \text{signes différents}\Rightarrow -}.
\]

L'inverse d'un nombre non nul $a$ est :
\[
\boxed{\dfrac1a}.
\]

Pour les fractions :
\[
\boxed{\dfrac ab\times\dfrac cd=\dfrac{ac}{bd}}
\]
et :
\[
\boxed{\dfrac ab\div\dfrac cd=\dfrac ab\times\dfrac dc}.
\]

Un nombre rationnel peut avoir plusieurs écritures ; choisir et contrôler l'écriture adaptée fait partie du raisonnement.
\end{aretenir}
